\section{Matching channel-separable order for the coupled family}\label{sec:matching-modes}
The obstruction in Appendix~\ref{sec:coupled} is sharp in order for its
precise family and row class. We construct a nonlocal, channel-separable
map from real coexact Laplacian modes. This supplies no general-profile
upper bound and no rectangularity claim for these new rows.

Let $r_{\min}$ be the minimum actual restricted rank among finite
channel-separable linear maps on $\mathcal U$ for which some penalty
vanishing at zero restores $c$-coercivity for the fixed profile
\eqref{eq:coupled-profile}. This class is nonempty: observe a complete
basis of $\mathcal U$, whose squared observations with coefficient $c$
control its norm since the energy is positive. Possible ranks are
integers between zero and $\dim\mathcal U$, so the minimum rank exists.
Define a separate finite Green-trace constant
\begin{equation}\label{eq:green-constant}
 \mathfrak G_D=D\,3^{D-1}/8.
\end{equation}
It is distinct from the scalar weighted-gap constant $G_D$ in
\eqref{eq:constants}.

\begin{theorem}[Matching order in the fixed-basis class]\label{thm:matching-modes}
For $D\ge3$, $n\ge2$, $g=V=n^D$ and the exact profile
\eqref{eq:coupled-profile},
\begin{equation}\label{eq:matching-order}
 \frac{D-1}{4(1+8D)}n^{2D-2}
 \le r_{\min}\le (D-1)(1+4\mathfrak G_D)n^{2D-2}.
\end{equation}
The upper map constructed below has a finite quadratic penalty with
$K=c+(4DH)^2/c$. Thus $r_{\min}=\Theta_D(n^{2D-2})$ for this family
and observation class. This is not an exact leading constant, efficient
algorithm, all-profile optimum or uniform-penalty result.
\end{theorem}

\subsection{Real modes and the exact energy decomposition}
The scalar forward torus Laplacian
$\mathscr L=\sum_i(D_i^+)^*D_i^+$ acts componentwise on spatial vectors.
It is real symmetric and commutes with curl and its adjoint, hence
preserves $\mathcal U_0$. Its eigenvalue at nonzero frequency $k$ is
\[
 \lambda(k)=4\sum_i\sin^2(\pi k_i/n)>0.
\]
There are $D-1$ coexact polarizations at every nonzero frequency and
none at zero, by \eqref{eq:closed-decomposition}. Conjugate pairs give
$2(D-1)$ real modes for those two frequencies, and self-conjugate
frequencies give $D-1$ real modes. Equivalently the real symmetric
restriction has a real counting-orthonormal eigenbasis
$f_s$, $1\le s\le d_0=(D-1)(V-1)$, with this same full eigenvalue
multiset. There is no further real/complex factor of two.
The torus estimate \eqref{eq:torus-gap} and the symbols give
\[
 16/n^2\le\lambda_s\le4D.
\]
Expand $z_\alpha=\sum_sz_{\alpha,s}f_s$ and put
$z_s=(z_{\alpha,s})_{\alpha=1}^g$. Orthonormality gives
\begin{equation}\label{eq:mode-energy}
 \|z\|^2=\sum_s\|z_s\|^2,\qquad
 Q_A(z)=\sum_s\lambda_s z_s^{\mathsf T}Bz_s.
\end{equation}
For the second identity,
$\langle df_s,df_t\rangle=\langle f_s,\mathscr Lf_t\rangle
=\lambda_t\delta_{st}$; the constant internal matrix $B$ does not
alter that calculation. It also holds within degenerate eigenspaces.
No independence of the individual spatial components as coexact domains
is asserted.

\subsection{Selected real coordinates and a strict kernel margin}
For each real mode let
\begin{equation}\label{eq:mode-selection}
 k_s=\min\left\{g,
 \left\lceil\max\left\{0,\frac{g(2c/\lambda_s-a)}{H-a}\right\}\right\rceil
 \right\}.
\end{equation}
Observe $V^{-1/2}z_{\alpha,s}$ for $\alpha\le k_s$.
Each row uses one internal channel and may mix sites and spatial components.
The row representers $e_\alpha f_s$ are an orthonormal subset on
$\mathcal U$, so the actual restricted rank is exactly
\begin{equation}\label{eq:selected-rank}
 r_C=\sum_sk_s.
\end{equation}
The rows are selected from the fixed profile/mode data before the field.
They extend to all fields by the same inner products, but the inequality
claimed here is on $\mathcal U$.

In a kernel field, the first $k_s$ internal coordinates of mode $s$
vanish. If $k_s=g$ nothing remains. Otherwise the remaining principal
block of $B$ is
\[
 H I_{g-k_s}-\frac{H-a}{g}\,11^{\mathsf T},
\]
whose least eigenvalue is $a+(H-a)k_s/g$. The ceiling in
\eqref{eq:mode-selection} implies
$\lambda_s[a+(H-a)k_s/g]\ge2c$ whenever its positive argument applies.
If the maximum argument is zero, $a\lambda_s\ge2c$ already.
If the outer truncation applies there is no unobserved part of that mode.
Thus \eqref{eq:mode-energy} proves the strict kernel bound
\begin{equation}\label{eq:kernel-margin}
 Q_A(x)\ge2c\|x\|^2\quad(x\in\ker(C|_{\mathcal U})).
\end{equation}
These are deliberately chosen whole real Laplacian modes; no commutation
of an arbitrary observation projection with the energy was assumed.

\subsection{A finite penalty on all of the actual domain}
Let $R$ be the counting orthogonal projection onto the selected row space
inside $\mathcal U$, and split $z=x+y$, $x=(I-R)z$, $y=Rz$.
Then $x$ is in the kernel, $x\perp y$, and
$\|Cz\|^2=V^{-1}\|y\|^2$ exactly.
The positive energy operator on $\mathcal U$ has norm at most
$M_*=4DH$, from $A\le HI$ and the full gradient norm bound $4D$.
Dropping $Q_A(y)\ge0$, applying \eqref{eq:kernel-margin} and bounding
the cross term gives
\begin{align*}
 Q_A(x+y)&\ge2c\|x\|^2-2M_*\|x\|\|y\|\\
 &\ge c\|x\|^2-(M_*^2/c)\|y\|^2.
\end{align*}
The last step is Young's inequality. Hence the finite choice
$K=c+M_*^2/c$ yields
\[
 V^{-1}Q_A(z)+K\|Cz\|^2\ge c\normn{z}^2.
\]
The margin is $2c$, not merely $c$: one unit of $c$ was spent controlling
the cross term. No uniform or optimal penalty is inferred.

\subsection{Ceiling count and the finite Green trace}
Since $a/H=1/(16Dn^2)<1/2$, $H-a\ge H/2$.
Whether the maximum argument is positive, zero or truncated,
\[
 k_s\le1+\frac{4g}{n^2\lambda_s},\qquad
 r_C\le d_0+\frac{4g}{n^2}\sum_s\lambda_s^{-1}.
\]
The ceiling cost is $d_0$, rather than $gd_0$.
Choose centered integer representatives of each frequency with
$|k_i|\le n/2$, taking only one Nyquist representative in each even
coordinate. The sine chord estimate gives
$\lambda(k)\ge16|k|_2^2/n^2$.
For $N=\lfloor n/2\rfloor$, overcount nonzero frequencies by every
integer point on max-norm shells $\ell=1,\ldots,N$.
Each shell has at most
\[
 (2\ell+1)^D-(2\ell-1)^D
 \le2D(2\ell+1)^{D-1}
 \le2D\,3^{D-1}\ell^{D-1}
\]
points. The first inequality follows by integrating $Dt^{D-1}$ over
the interval of length two. On that shell $|k|_2\ge\ell$.
Since $D\ge3$,
$\sum_{\ell=1}^N\ell^{D-3}\le N^{D-2}\le n^{D-2}$,
including $D=3$. Therefore the complete finite sum satisfies
\begin{equation}\label{eq:finite-green}
 \sum_{k\ne0}\lambda(k)^{-1}
 \le\frac{n^2}{16}\,2D3^{D-1}\sum_{\ell=1}^N\ell^{D-3}
 \le\mathfrak G_DV.
\end{equation}
Odd sizes and Nyquist overcounting only enlarge this upper estimate.
Multiplying by the real coexact multiplicity gives
$\sum_s\lambda_s^{-1}\le(D-1)\mathfrak G_DV$. Thus
\[
 r_C\le(D-1)(V-1)+4(D-1)\mathfrak G_DgV/n^2
 \le(D-1)(1+4\mathfrak G_D)n^{2D-2}.
\]
Here $g=V=n^D$ and $V\le n^{2D-2}$. The estimate need not beat the
trivial full dimension at small sizes; the actual truncation still
ensures $r_C\le gd_0$.

Finally Appendix~\ref{sec:coupled} applies to every allowed map and any
penalty vanishing at zero. With $V-1\ge V/2$ and
$1+8Dn^2\le(1+8D)n^2$, its necessary bound is at least
$[(D-1)/(4(1+8D))]n^{2D-2}$. Combining proves
\eqref{eq:matching-order}. The mixed map of the preceding appendix still
has actual rank $d_0=O_D(V)$ on this same profile. Matching the separable
order therefore strengthens the fixed-basis comparison without producing
an all-linear, basis-free, fixed-$g$, all-profile or physical obstruction.
The Fourier, principal-block, Young and finite-shell mechanisms used here
are classical; no historical-priority or formal-execution claim is made.
